\chapter{Run, Monitor, and Diagnose Remote Work}\label{ch:remote}
\section{Preflight the actual execution boundary}

Use this flow when validation needs data, software, or resources outside the current
writing environment, or when a job outlives an interactive step. Start with the owning
repository's host and execution documentation. Establish the literal host, exact target
and owner, capability, and consequence. Synthetic checks can run on the writing host;
service-dependent checks need the documented capable environment \cite[Remote execution
  and interruption]{skill:remote}.

Access to remote data does not by itself require remote execution. Choose the execution
host for its runtime, input access and capacity, then test the exact input from that
host. Before relying on a route, distinguish a verified read from an unverified
submission path: a successful storage read does not validate batch submission.
\hook{remote-hosts}

Probe the chosen SSH alias non-interactively before dependent work. With a shared control
master, issue one bounded command at a time through that alias. After a timeout, broken
pipe, multiplexing error, or unexpected disconnect, stop issuing calls and inspect the
master once. Do not replace the user's transport as an incidental repair. If
authentication is absent, the user establishes one interactive login; the agent never
requests a password or certificate passphrase \cite[SSH control masters]{skill:remote}.

The authorization must name the intended remote action. Inspecting logs is different from
submitting or cancelling a job. A bounded pilot may be a useful first execution when its
submission is authorized; use it to check the input, configuration, output schema, and
actual resource behavior before expanding the run. Keep unbounded analysis inside the
scheduler rather than on a login node. Source edits stay in the explicitly owned local
checkout or worktree, never through a mounted remote filesystem \cite[Remote checkouts
  and execution; Scheduler jobs]{skill:remote}.

\section{Preserve the parent contract before waiting}

Record exact job IDs, status surfaces, evidence paths, allowed and forbidden actions,
monitoring interval and deadline, terminal validation, and downstream artifacts. Include
separately any report, review, source publication, delivery, and notification
requirements. This lets the parent recover the entire task when a wait returns rather
than treating scheduler completion as the endpoint \cite[Durable completion
  contract]{skill:jobs}.

Prefer an existing authorized programmatic monitor. It can observe exact external tasks
and write local snapshots without spending model calls on polling. The documented monitor
preserves atomic current state, append-only observations, a terminal result, and an
exclusive lock. Snapshot expiry matters: a stale record is missing current evidence, not
proof of a successful task \cite[Monitor selection; Programmatic job
  monitors]{skill:jobs}.

Where the client provides one, a parent wait call binds a reviewed contract and monitor
hashes, a finite duration, a handback policy, and any exact notification authorization.
It returns evidence to the parent. A client that lacks such a tool must verify its own
integration and parent identity handling before relying on that behavior. Otherwise
record the gap and use the documented bounded fallback or explicit manual recovery of
existing evidence \cite[Client scope; Parent wait contract]{skill:jobs}.

Before unattended work, preflight the exact external read and write roots, permission
configuration, and ownership. A permission prompt is a handback condition; widening paths
or using blanket auto-approval does not solve the contract. The specific configuration
rules belong to each client, while the requirement to respect the task boundary applies
across clients \cite[Permission and access preflight; Client scope]{skill:jobs}.

\section{Interpret the handback before changing anything}

A monitor supplies observations. Table~\ref{tab:handback} describes the next reasoning
step for common handbacks; repair or retry requires its own recorded authorization.

\begingroup
\small
\begin{longtable}{L{0.27\textwidth}L{0.66\textwidth}}
  \caption{Remote handback and the bounded next inspection
    \cite[Programmatic job monitors; Terminal handling]{skill:jobs}.}
  \label{tab:handback}\\
  \toprule Observation & Next inspection \\ \midrule
  \endfirsthead
  \toprule Observation & Next inspection \\ \midrule
  \endhead
  Active, fresh snapshot & Continue the contracted wait within its lease. \\
  Stale or absent snapshot & Establish monitor liveness and current evidence;
  do not infer external job success. \\
  Terminal clean & Validate the predeclared files, events, schema, checksums,
  provenance, normalization, and storage before downstream use. \\
  Terminal with failures & Inspect the exact failing task/index and relevant
  logs; preserve successful outputs and failed-attempt evidence. \\
  Owner/scope/invariant drift & Reconcile the exact job identities, specification,
  and expected counts; stop expansion until the discrepancy is understood. \\
  Transport failure & Apply the shared SSH stop-and-inspect rule; distinguish
  loss of observation from failure of the analysis. \\
  \bottomrule
\end{longtable}
\endgroup

Suppose an array reports mixed completion. First identify the missing or failed indices
and their exit evidence, then check which outputs exist and which input coverage they
represent. Preserve successful evidence. The smallest next decision may concern only the
failed subset, but monitoring permission alone does not authorize its resubmission.
Record the cause, bounded proposal, and required authorization instead of restarting the
entire production by assumption \cite[Terminal handling]{skill:jobs}.

If the scheduler reports success but event coverage is incomplete, downstream analysis
remains held. Follow the upstream validation order in Chapter~\ref{ch:code}: input and
accounting before distributions and final results. Logs, schema inspection, and exact job
evidence should drive the diagnosis. Add focused logging only when the current evidence
cannot distinguish the plausible causes \cite[Validation order]{skill:implementation}.

\section{Handle client loss and notification uncertainty}

A documented wait separates observation state from delivery state. A detached notifier
can continue observing existing monitors after client loss and record an orphaned result;
it never starts an agent or model. An explicitly selected recovery owner must reconcile
the result and the remaining contract \cite[Client loss and notification; Compaction and
  recovery]{skill:jobs}.

Email is disabled unless the exact command and destination are authorized. An interrupted
attempt can leave delivery externally uncertain. Do not retry an ambiguous attempt
automatically: reconcile the notifier and transport evidence first. Acknowledging a wait
means its evidence was consumed; it does not validate the output or complete the
scientific task \cite[Client loss and notification]{skill:jobs}.

At handback, update the parent state with terminal checks, known failures, remaining
analysis or report work, and one next safe action. Restore the same obligations after
compaction. A job leaving the queue, a worker returning, a PDF compiling, and a local
email command succeeding each support a particular claim; none alone closes the whole
parent contract \cite[Completion evidence]{skill:jobs}.
